%===========================================================
%  MediSeg — User Guide
%  Smart Mobile Medical-Waste Collection & Segregation System
%  Autodesk · Smart India Hackathon 2026
%
%  Compile:  pdflatex UserGuide.tex
%===========================================================
\documentclass[11pt, a4paper]{article}

%------------- Packages --------------
\usepackage[margin=2.4cm]{geometry}
\usepackage[utf8]{inputenc}
\usepackage[T1]{fontenc}
\usepackage{amsmath, amssymb}
\usepackage{graphicx}
\usepackage{xcolor}
\usepackage[colorlinks=true, linkcolor=teal!75!black, urlcolor=teal!75!black]{hyperref}
\usepackage{enumitem}
\usepackage{booktabs}
\usepackage{tabularx}
\usepackage{listings}
\usepackage{fancyhdr}
\usepackage[most]{tcolorbox}
\usepackage{titlesec}

%------------- Colours --------------
\definecolor{mediseg}{HTML}{14b8a6}      % teal
\definecolor{medisegdark}{HTML}{0f766e}  % deep teal
\definecolor{slate}{HTML}{334155}
\definecolor{codebg}{HTML}{f1f5f9}

%------------- Page style --------------
\pagestyle{fancy}
\fancyhf{}
\fancyhead[L]{\small\textcolor{mediseg}{\textbf{MediSeg -- User Guide}}}
\fancyhead[R]{\small\textcolor{slate}{SIH 2026 \textperiodcentered\ Autodesk}}
\fancyfoot[C]{\small\thepage}
\renewcommand{\headrulewidth}{0.4pt}
\headrulecolor{mediseg!60!white}

%------------- Section titles --------------
\titleformat{\section}
  {\Large\bfseries\color{medisegdark}}
  {\thesection}{0.8em}{}
\titleformat{\subsection}
  {\large\bfseries\color{slate}}
  {\thesubsection}{0.7em}{}
\titleformat{\subsubsection}
  {\normalsize\bfseries}
  {\thesubsubsection}{0.6em}{}

%------------- Listing style --------------
\lstdefinestyle{cmd}{
  language=bash,
  backgroundcolor=\color{codebg},
  basicstyle=\ttfamily\small,
  keywordstyle=\color{medisegdark}\bfseries,
  commentstyle=\color{gray},
  showstringspaces=false,
  frame=single,
  rulecolor=\color{mediseg!30!white},
  breaklines=true,
  columns=fullflexible
}

%------------- Boxes --------------
\newtcolorbox{tipbox}{
  colback=teal!8!white, colframe=mediseg,
  title=\textbf{Note}, fonttitle=\bfseries, arc=3pt, boxrule=0.8pt
}
\newtcolorbox{warnbox}{
  colback=orange!8!white, colframe=orange!85!black,
  title=\textbf{Important}, fonttitle=\bfseries, arc=3pt, boxrule=0.8pt
}

%===========================================================
\begin{document}

%------------- Title --------------
\begin{titlepage}
\centering
\vspace*{2.5cm}

{\color{teal!50!white}\rule{\textwidth}{1pt}}\\[0.8em]
{\Huge\bfseries\textcolor{mediseg}{MediSeg}}\\[0.5em]
{\LARGE Smart Mobile Medical-Waste\\[0.3em] Collection \& Segregation System}\\[1.2em]
{\color{teal!50!white}\rule{\textwidth}{1pt}}\\[2.5em]

{\large \textbf{User Guide}}\\[1.2em]
{\large Web Application \textperiodcentered\ Interactive Prototype}\\[3em]

\begin{minipage}{0.75\textwidth}
\centering
\small
Autodesk \textperiodcentered\ Smart India Hackathon 2026\\
Repo: \url{https://github.com/Manjunath2006548/mediseg-sih2026}\\
Live: \url{https://mediseg-sih2026.vercel.app}
\end{minipage}

\vfill
{\small Version 1.0 \textperiodcentered\ September 2026}
\end{titlepage}

%------------- TOC --------------
\tableofcontents
\clearpage

%===========================================================
\section{Introduction}

\textbf{MediSeg} is a web-based solution for the Autodesk problem statement
\textit{``Design a Smart Mobile Medical-Waste Collection and Segregation System''}.
Hospital bio-medical waste is hazardous to handle; when mis-segregated it is
either incinerated incorrectly or leaks into the environment. MediSeg replaces
manual pushcarts with an \emph{autonomous mobile rover} that:

\begin{itemize}[leftmargin=1.6em, itemsep=2pt]
  \item navigates hospital zones autonomously (LiDAR-based patrol),
  \item recognises waste bags with an on-board AI camera (YOLOv8-style model, ~98\,\%),
  \item drops each bag into the correct BMWM 2018 colour-coded compartment, and
  \item records a full QR chain-of-custody for regulatory compliance.
\end{itemize}

This guide walks through the delivered web prototype: the login flow, the four
pages of the application, the interactive simulation, and deployment.

\subsection{What is provided}
\begin{tabularx}{\textwidth}{lX}
  \toprule
  \textbf{Item} & \textbf{Detail}\\
  \midrule
  Web application & Next.js static export (no backend required)\\
  Interactive prototype & Autonomous patrol \& AI scanner simulator\\
  Concept deck & Problem-solution presentation slides\\
  GitHub repository & \url{https://github.com/Manjunath2006548/mediseg-sih2026}\\
  Live deployment & \url{https://mediseg-sih2026.vercel.app}\\
  \bottomrule
\end{tabularx}

%===========================================================
\section{System Requirements}

\begin{tabularx}{\textwidth}{lX}
  \toprule
  \textbf{Category} & \textbf{Requirement}\\
  \midrule
  Operating system & Windows / macOS / Linux (any modern browser)\\
  Browser & Chrome 100+, Edge, Firefox, Safari (latest)\\
  Screen & Desktop recommended for full prototype experience\\
  JavaScript & Enabled (the app runs entirely in the browser)\\
  Internet & Needed only for first page load (Google Fonts) and hosting\\
  Node.js (developers) & Node 18+ to build locally ($\ge$ 18.17 recommended)\\
  \bottomrule
\end{tabularx}

\begin{tipbox}
The prototype runs fully client-side. No backend, database, or API key is
required, so it works offline once the page is loaded.
\end{tipbox}

%===========================================================
\section{Getting Started (End User)}
\label{sec:getstarted}

\subsection{Accessing the application}
\begin{enumerate}[leftmargin=1.6em, itemsep=2pt]
  \item Open \url{https://mediseg-sih2026.vercel.app} in a browser.
  \item The \textbf{portal page} loads and shows the project overview plus four
        cards: \textit{Introduction}, \textit{Sign In / Register},
        \textit{Working Prototype}, and \textit{Concept Deck}.
  \item Click any card, or use \textbf{Launch Simulation} to go straight to the
        prototype.
\end{enumerate}

\subsection{Registration}
\begin{enumerate}[leftmargin=1.6em, itemsep=2pt]
  \item Navigate to the \textbf{Sign In / Register} card.
  \item Switch to the \textbf{Register} tab.
  \item Enter a name, a valid email, and a password. The password must satisfy
        the \textbf{strong-password policy}:
\end{enumerate}
\begin{itemize}[leftmargin=1.6em, itemsep=2pt]
  \item at least 8 characters,
  \item at least one uppercase letter,
  \item at least one lowercase letter,
  \item at least one digit,
  \item at least one special character.
\end{itemize}
A live \textbf{strength meter} and a per-rule checklist show progress as you
type. Weak passwords are rejected with a clear message.

\subsection{Signing in}
\begin{enumerate}[leftmargin=1.6em, itemsep=2pt]
  \item Enter the email and password you registered with (any valid credentials
        are accepted in this demo).
  \item On success you are taken to the prototype with a personalised header.
  \item Your session is stored locally in the browser and persists across
        reloads.
\end{enumerate}

\begin{warnbox}
This is a prototype, so credentials are validated in the browser only.
Replace the client-side check with a real authentication service for
production use.
\end{warnbox}

%===========================================================
\section{Application Pages}

\begin{tabularx}{\textwidth}{lX}
  \toprule
  \textbf{Route} & \textbf{Purpose}\\
  \midrule
  |/| & High-level portal linking all sections\\
  |/mediseg/intro| & Problem background, impact, architecture\\
  |/mediseg/login| & Login / registration with strong-password check\\
  |/mediseg| & \textbf{Main interactive prototype} (login required)\\
  |/mediseg/presentation| & Concept deck (presentation slides)\\
  \bottomrule
\end{tabularx}

\subsection{Portal (|/|)}
Default landing page. Contains the project tagline, module summaries
(A: Navigation \& Collection, B: AI Recognition \& Segregation) and the
all-in-one portal cards that link to every part of the demonstration.

\subsection{Introduction (|/mediseg/intro|)}
Explains the problem statement: the hazards of improper medical-waste handling,
the biomedical-waste rules, why manual segregation is unsafe, and how the
rover-based pipeline works end to end.

\subsection{Authentication (|/mediseg/login|)}
Handles login and registration as described in
Section~\ref{sec:getstarted}.

\subsection{Concept Deck (|/mediseg/presentation|)}
A slide-style presentation covering problem context, proposed solution,
system design, methodology, and expected outcomes. Use the section anchors or
page down through the deck.

%===========================================================
\section{The Interactive Prototype}

The core of the demonstration is the prototype at \url{/mediseg}. It simulates
the complete collection-and-segregation workflow in four linked panels.

\subsection{Robot View \textperiodcentered\ Autonomous Patrol}
\begin{itemize}[leftmargin=1.6em, itemsep=2pt]
  \item An \textbf{SVG floor plan} represents the hospital layout with a patrol
        route and waste drop points.
  \item The rover \textbf{moves along the patrol path} automatically, stopping
        at each depot to collect a generated waste bag.
  \item \textbf{Obstacle avoidance} is simulated -- the rover steers around
        obstacles on the map.
  \item Use the \textbf{speed slider} to slow down or speed up the simulation,
        and toggle \textbf{auto-run} to pause/resume.
\end{itemize}

\subsection{AI Scanner \textperiodcentered\ Waste Recognition}
\begin{itemize}[leftmargin=1.6em, itemsep=2pt]
  \item A simulated \textbf{camera feed} displays the bag currently being
        inspected, overlaid with a detection bounding box.
  \item The model reports the \textbf{category} (e.g.\ yellow bag), a
        \textbf{confidence score}, and a simulated \textbf{accuracy of 98\,\%}.
  \item When no bag is present the camera shows an \textbf{idle state}
        (``Camera Live -- Awaiting Waste Intake'') with a
        \textbf{Generate Waste Now} button to spawn a new bag.
\end{itemize}

\subsection{Segregation \textperiodcentered\ BMWM 2018 Compartments}
\begin{itemize}[leftmargin=1.6em, itemsep=2pt]
  \item Five \textbf{colour-coded compartments} match the Bio-Medical Waste
        (Management)\ Rules 2018 categories.
  \item The sorting mechanism delivers each recognised bag into the correct
        compartment automatically.
  \item Fill levels animate as bags are dropped, giving a clear visual summary.
\end{itemize}

\subsection{Tracking Ledger \textperiodcentered\ Compliance Report}
\begin{itemize}[leftmargin=1.6em, itemsep=2pt]
  \item Every recognised bag receives a \textbf{QR-style record ID}.
  \item The \textbf{ledger table} logs: timestamp, bag ID, detected category,
        assigned compartment, confidence, and collection status.
  \item A real-time \textbf{compliance report} computes the percentage of bags
        segregated correctly and highlights any mis-segregation.
\end{itemize}

\subsection{Controls and Session}
\begin{itemize}[leftmargin=1.6em, itemsep=2pt]
  \item \textbf{Battery indicator} simulates charge level and recharging cycles.
  \item \textbf{Logout} clears the local session and returns to the login page.
  \item The header shows the currently signed-in user.
\end{itemize}

%===========================================================
\section{Running a Demo Session (Scenario)}

Follow this script for a clean 2-minute demo:

\begin{enumerate}[leftmargin=1.6em, itemsep=3pt]
  \item Login to the prototype (Section~\ref{sec:getstarted}).
  \item On \textbf{Robot View}, press \emph{auto-run} -- the rover departs on
        its patrol route and reaches the first depot.
  \item Point out the \textbf{camera feed}: a bag appears, the model draws a
        bounding box and announces the category and confidence.
  \item Watch the bag slide into the correct \textbf{coloured compartment}
        and the compartment fill level rise.
  \item Open the \textbf{tracking ledger} and show a generated record (bag ID,
        category, compartment).
  \item Show the \textbf{compliance report} -- with correct segregation the
        score is 100\,\%; drop a mis-sorted bag to demonstrate the warning.
  \item Mention the \textbf{speed slider} as a stand-in for real-time telemetry.
\end{enumerate}

%===========================================================
\section{Developer Guide}

\subsection{Folder layout}
\begin{lstlisting}[style=cmd]
mediseg-sih2026/
|-- app/
|   |-- page.js                  # portal (landing)
|   |-- layout.js                # root layout + metadata
|   |-- globals.css              # Tailwind + custom utilities
|   `-- mediseg/
|       |-- page.js              # interactive prototype
|       |-- intro/page.js        # introduction page
|       |-- login/page.js        # login / register
|       `-- presentation/page.js # concept deck
|-- public/
|   `-- images/mediseg-bg.svg    # themed background
|-- next.config.js               # static export config
|-- tailwind.config.js
`-- package.json
\end{lstlisting}

\subsection{Local development}
\begin{lstlisting}[style=cmd]
npm install      # install dependencies
npm run dev      # start dev server on http://localhost:3000
\end{lstlisting}

\subsection{Production build}
\begin{lstlisting}[style=cmd]
npm run build    # static export written to ./out
\end{lstlisting}

\begin{tipbox}
The pipeline is fully static (\texttt{output: 'export'}). The generated
\texttt{out/} directory can be hosted by any static file server.
\end{tipbox}

\subsection{Deploy to Vercel}
\begin{lstlisting}[style=cmd]
vercel --prod --yes
\end{lstlisting}
After the first deployment the project is linked, so subsequent deploys are
just \texttt{vercel --prod}. Alternatively, connect the GitHub repository and
enable auto-deploy on push.

\subsection{Deploy to any static host}
Since the build output is plain HTML/CSS/JS:
\begin{enumerate}[leftmargin=1.6em, itemsep=2pt]
  \item run \texttt{npm run build},
  \item upload the contents of \texttt{out/} to Netlify, GitHub Pages,
        or an Nginx/S3 bucket.
\end{enumerate}

%===========================================================
\section{Troubleshooting}

\begin{tabularx}{\textwidth}{lX}
  \toprule
  \textbf{Symptom} & \textbf{Resolution}\\
  \midrule
  Page shows stale version & Hard-refresh (Ctrl+F5 / Cmd+Shift+R).\\
  Login does not persist & Your browser blocks \texttt{localStorage}; allow
  site data or try another browser.\\
  Prototype ``locked'' & You are not signed in. Register or sign in, then open
  the prototype again.\\
  Robot not moving & Ensure \emph{auto-run} is enabled and raise the speed
  slider.\\
  Scanner shows idle & Expected idling state; press \textbf{Generate Waste Now}
  to inject a bag.\\
  Weak-password error & Password violates one of the five rules
  (Section~\ref{sec:getstarted}); use the checklist.\\
  Deployment not updating & Run \texttt{vercel --prod --yes} from the
  repository root after rebuilding.\\
  \bottomrule
\end{tabularx}

%===========================================================
\section{Credits}

\begin{itemize}[leftmargin=1.6em, itemsep=2pt]
  \item \textbf{Team:} Manjunath Hebbar and collaborators
        (\texttt{Manjunath2006548} on GitHub).
  \item \textbf{Organisers:} Smart India Hackathon 2026 \textperiodcentered\
        Autodesk.
  \item \textbf{Regulatory basis:} Bio-Medical Waste (Management) Rules, 2018.
  \item \textbf{Stack:} Next.js 14, React 18, Tailwind CSS, Vercel.
\end{itemize}

\end{document}